\begin{enumerate}[label=\textbf{\ansnum.\arabic*}]
  \item (a) Agent: the controller. State: joint angles, body tilt, contact of hands and feet. Action: target angles. Reward: positive for height of the body or for standing, negative for falling back. (b) State: position of ball, robot and goal. Action: walking command. Reward: positive when the ball gets closer to the goal, large reward for a goal. (c) Agent: the thermostat. State: room temperature. Action: heater on or off. Reward: negative for the difference to the wanted temperature (and for energy use).
  \item $2(1+0.9+0.81+0.729)=2\cdot3.439=6.878$.
  \item $0+0.5\cdot0+0.25\cdot5=1.25$.
  \item $\gamma^3\cdot10=8\Rightarrow\gamma^3=0.8\Rightarrow\gamma=0.8^{1/3}\approx0.928$.
  \item All ten rewards are $-1$: $G=-\sum_{k=0}^{9}0.9^k=-\frac{1-0.9^{10}}{0.1}\approx-6.51$.
  \item Almost all episodes of a robot that has not learned anything end with reward $0$, so there is nothing to learn from (the reward is \emph{sparse}). An additional reward for getting closer to the ball, or for staying upright, gives feedback in every step.
  \item $Q_1=0$, $Q_2=0+\frac12(1-0)=0.5$, $Q_3=0.5+\frac13(1-0.5)=0.667$, $Q_4=0.667+\frac14(0-0.667)=0.5$, $Q_5=0.5+\frac15(1-0.5)=0.6$. (These are the averages $0,\frac12,\frac23,\frac24,\frac35$.)
  \item Greedy action: $0.9+\frac{0.1}{3}\approx0.933$. Each other action: $\frac{0.1}{3}\approx0.033$.
  \item Example:
\begin{verbatim}
import random
p = [0.3, 0.5, 0.8]
Q = [0.0, 0.0, 0.0]
N = [0, 0, 0]
eps = 0.1
for t in range(1000):
    if random.random() < eps:
        a = random.randrange(3)
    else:
        a = max(range(3), key=lambda i: Q[i])
    r = 1 if random.random() < p[a] else 0
    N[a] += 1
    Q[a] += (r - Q[a]) / N[a]
print(Q, N)
\end{verbatim}
  The third kick (true rate $0.8$) is used most often, and $Q_C$ is close to $0.8$. The estimates of A and B are less accurate because they are tried rarely.
  \item[\textbf{\ansnum.10}] $G=1+\gamma+\gamma^2+\dots=\frac{1}{1-\gamma}$ (geometric series). For $\gamma=0.99$: $\frac{1}{0.01}=100$.
\end{enumerate}
